\documentclass[letterpaper,11pt]{article}

% --- PACKAGES ---
\usepackage[T1]{fontenc}
\usepackage{xcolor}
\usepackage{geometry}
\usepackage{parskip}
\usepackage{enumitem}
\usepackage{titlesec}
\usepackage{fontawesome5}
\usepackage{hyperref}
\usepackage{fontspec}
\usepackage{paracol} % For two-column layout

% --- DOCUMENT SETUP ---
\geometry{
    left=0.5in,
    right=0.5in,
    top=0.6in,
    bottom=0.6in
}

% --- FONT & COLOR DEFINITIONS ---
\definecolor{primarycolor}{HTML}{B22222}
\definecolor{graytext}{HTML}{555555}
\definecolor{linkcolor}{HTML}{005F6B}

\setmainfont{Lato}
\setsansfont{Lato}

% --- HYPERLINK SETUP ---
\hypersetup{
    colorlinks=true,
    urlcolor=linkcolor,
    linkcolor=linkcolor,
    citecolor=linkcolor,
}

% --- HEADER & FOOTER ---
\usepackage{fancyhdr}
\pagestyle{fancy}
\fancyhf{}
\renewcommand{\headrulewidth}{0pt}
\renewcommand{\footrulewidth}{0pt}
\fancyfoot[R]{\small\thepage}

% --- CUSTOM COMMANDS & FORMATTING ---
\setlength{\parindent}{0pt}

\setlist[itemize]{
    itemsep=2.5pt,
    parsep=0pt,
    topsep=0pt,
    leftmargin=1.2em
}

\setlength{\tabcolsep}{1pt}

\titleformat{\section}
    {\Large\bfseries\color{primarycolor}}
    {}
    {0em}
    {}[\color{primarycolor}\titlerule]
\titlespacing*{\section}{0pt}{8pt}{4pt}

\newcommand{\entry}[4]{%
    \textbf{#1} \hfill \textit{\color{graytext} #2} \\
    \textit{#3} \hfill \textit{\color{graytext} #4}\par
}

% --- RESUME CONTENT ---
\begin{document}

{\raggedright
{\Huge\bfseries\color{primarycolor}{Arun Jayaraman}} \\[4pt]
\faMapMarker* Rockville, MD $\vert$ \faPhone* (847) 436-3386 $\vert$ \faEnvelope \ \href{mailto:arun160638@gmail.com}{arun160638@gmail.com} $\vert$ \faLinkedin \ \href{https://linkedin.com/in/arun-jayaraman}{/in/arun-jayaraman} $\vert$ \faGraduationCap \ \href{https://scholar.google.com/citations?user=9L9vIUcAAAAJ}{Google Scholar}
\par}

% --- Summary ---
\section*{Summary}
Senior Scientist and translational data scientist with 6+ years of applied computational research across pharmaceutical R\&D and academia, specializing in immunology, single-cell and multimodal omics, biomarker strategy, and precision medicine. Leading computational strategy for RA clinical development and cross-indication programs, translating complex molecular and clinical data into actionable biomarker, investment, and portfolio decisions.

% --- Work Experience ---
\section*{Work Experience}
\entry{Senior Scientist}{Jan 2023 -- Present}{AstraZeneca}{Gaithersburg, MD}
\begin{itemize}
    \item Primary translational data scientist for AZ's early and clinical-stage RA programs, leading biomarker strategy and multi-omics analyses supporting translational, clinical development, and investment decisions.
    \item Lead computational strategy across academic and consortium collaborations spanning pre-RA through established autoimmune disease, integrating longitudinal multi-omic, single-cell, spatial, immune-repertoire, and clinical data to inform disease positioning and biomarker development.
    \item Develop advanced single-cell and multimodal analyses across internal and public datasets, including a scVI-based cross-disease T-cell atlas, matched blood/tissue studies, and integrated CITE-seq analyses.
    \item Influence portfolio strategy across RA, pediatric JIA, cell therapy, TCEs, and multispecific programs through indication assessments, biomarker analyses, competitive positioning, and synthesis of molecular and clinical evidence.
    \item Evaluate emerging biomarker technologies and guide assay selection, experimental design, cohort construction, and validation; benchmarked three autoantibody-discovery platforms and provided recommendations for further investment based on sensitivity and technical maturity.
    \item Co-lead development of AI-enabled workflows over internal omics repositories to combine validated internal data with external evidence and accelerate cross-indication assessment and indication prioritization.
    \item Served as project lead for a centralized R Shiny translational-data platform enabling self-service molecular and signature analyses; coordinated collaborative development and production deployment using GitHub and CI/CD.
\end{itemize}

\entry{Associate Scientist II}{Jan 2022 -- Jan 2023}{AbbVie Genomics Research Center}{North Chicago, IL}
\begin{itemize}
    \item Identified disease drivers in therapeutic non-responders through integrated single-cell and bulk transcriptomics, including a pathogenic synovial fibroblast population enriched in RA non-responders.
    \item Applied longitudinal mixed-effects modeling to upadacitinib-treated arthritis cohorts, validating drug-associated molecular signatures and informing therapeutic strategy.
\end{itemize}

\entry{Graduate Researcher}{Sept 2020 -- Dec 2021}{Emory University, The Davis Lab}{Atlanta, GA}
\begin{itemize}
    \item Led scRNA-seq studies of cardiac stromal/progenitor populations and collaborated with Yale School of Medicine on T-cell responses associated with cardiac allograft rejection, contributing to multiple peer-reviewed publications.
\end{itemize}

\entry{AI/ML-Driven Drug Discovery and Cheminformatics}{Summer 2020}{AbbVie Innovation Center}{Urbana, IL}
\begin{itemize}
  \item Applied generative models (VAE, GAN, RL) to \textit{de novo} molecular design and developed a hierarchical neural network to predict chemical reaction conditions for drug-discovery workflows.
\end{itemize}

\newpage

% ================= EDUCATION =================
\section*{Education}
\entry{Professional Certificate in Artificial Intelligence}{Expected Oct 2026}{Stanford School of Engineering, Artificial Intelligence Professional Program}{Online}
\entry{M.S., Biomedical Engineering}{Dec 2021}{Georgia Institute of Technology}{Atlanta, GA}
\entry{B.S., Chemical Engineering, With Distinction (Dean's List)}{May 2020}{University of Illinois Urbana-Champaign}{Urbana, IL}

% ================= SKILLS =================
\section*{Technical Skills}
\begin{itemize}
  \item \textbf{Single-cell \& multimodal:} scRNA-seq, CITE-seq, scVI, Seurat, Scanpy, Harmony, CCA, pseudobulk differential expression and abundance, Monocle~3, AUCell
  \item \textbf{Repertoire, cytometry \& spatial:} single-cell and bulk BCR/TCR-seq (Cell Ranger VDJ, scRepertoire); CyTOF; Xenium, GeoMx, Space Ranger
  \item \textbf{Bulk omics \& integration:} DESeq2, edgeR, limma, mixed-effects models (lme4, dream), MOFA+, WGCNA, GSVA, singscore; Olink proteomics; STAR, Salmon
  \item \textbf{Machine learning \& AI:} PyTorch, TensorFlow, scikit-learn;
deep generative models; reinforcement learning; AI agents;
AI coding assistants
  \item \textbf{Programming \& infrastructure:} Python, R, Bash, SQL, JavaScript; Git/GitHub, CI/CD; Linux HPC, AWS; Nextflow, Snakemake; R Shiny, ggplot2, Matplotlib
\end{itemize}

% ================= PUBLICATIONS =================
\section*{Selected Publications, Patents \& Presentations}
\begingroup\small

\textbf{Patent}
\begin{itemize}
  \item ``PAD ENZYMES AS BIOMARKERS'' (2025). U.S. Provisional Patent Application No.~63/811,164.
\end{itemize}

\textbf{Peer-reviewed publications} (* co-first author)
\begin{itemize}
  \item Korutla L, Hoffman JR, \ldots, \textbf{Jayaraman A}, et al. (2024).
  Circulating T cell specific extracellular vesicle profiles in cardiac
  allograft acute cellular rejection. \textit{Am. J. Transplant.}

  \item Hoffman JR*, \textbf{Jayaraman AR*}, Bheri S, Davis ME (2022).
  Single-cell RNA sequencing reveals distinct cardiac-derived stromal cell
  subpopulations. \textit{J. Cardiovasc. Dev. Dis.} 9:374.

  \item Gunasekaran M, Mishra R, Saha P, Morales D, Cheng W-C,
  \textbf{Jayaraman AR}, et al. (2022). Comparative efficacy and mechanism
  of action of cardiac progenitor cells after cardiac injury.
  \textit{iScience} 25:104656.

  \item Hoffman JR, Park H-J, Bheri S, \textbf{Jayaraman AR}, Davis ME
  (2022). Comparative computational RNA analysis of cardiac-derived
  progenitor cells and their extracellular vesicles.
  \textit{Genomics} 114:110349.

  \item Park J, \textbf{Jayaraman A}, Schrader AW, Hwang GW, Han H (2020).
  Controllable modulation of precursor reactivity for synthesis of
  high-quality quantum dots. \textit{Nat. Commun.}
\end{itemize}

\textbf{Conference presentations}
\begin{itemize}
  \item Csomor E, \textbf{Jayaraman A}, Stewart L, \ldots, McInnes IB, Goodyear CS (2025). Prevalence of IFN-high rheumatoid arthritis in early and established RA patients. EULAR 2025 (ABS0252). \textit{Ann. Rheum. Dis.} 84(Suppl.~1):1950--1951.
  \item Sharrack S, Yu X, \textbf{Jayaraman A}, \ldots, Mankia K (2025). Elevated serum peptidylarginine deiminase 4 (PAD4) levels in anti-CCP antibody positive at-risk individuals with arthralgia who progress to rheumatoid arthritis. Posters, EULAR 2025 (POS0067; \textit{Ann. Rheum. Dis.} 84(Suppl.~1)) and ACR Convergence 2025.
  \item Azim A, \textbf{Jayaraman A}, et al. (2024). Nasal sampling identifies disease-relevant signals in a real-world airways disease cohort. Oral presentation, ERS Congress 2024. \textit{Eur. Respir. J.} 64(Suppl.~68):OA3764.
  \item Faner R, Azim A, \textbf{Jayaraman A}, et al. (2024). The NOVELTY nasal transcriptome identifies mechanisms in asthma and COPD exacerbators. Poster, ERS Congress 2024. \textit{Eur. Respir. J.} 64(Suppl.~68).
  \item van den Berge M, \ldots, Azim A, \textbf{Jayaraman A}, et al. (2024). The nasal transcriptome reflects clinically relevant disease activity in NOVELTY, with a different signal in asthma and in COPD. Poster, ERS Congress 2024. \textit{Eur. Respir. J.} 64(Suppl.~68).
\end{itemize}

\par\endgroup
\end{document}